\chapter[(Джанака)]{Радість у самореалізації (Джанака)}

[2.1] Я справді чистий і в спокої. Свідомість поза природною причинністю. Довго я давав завдавати собі болю оманою.

[2.2] Як я даю світло цьому тілу, так даю його й всесвіту. Тож весь всесвіт --- мій, або, з іншого боку, ніщо не моє.

[2.3] Тепер, коли я відсторонився від тіла й речей, великою благодаттю Істина відкривається.

[2.4] Як хвилі, піна й бульбашки не відрізняються від води, так усе, що постає з Я, не відрізняється від Я.

[2.5] Пильно подивись на тканину --- і побачиш лише нитки. Пильно подивись на творіння --- і побачиш лише Я.

[2.6] Як цукор, добутий з цукрової тростини, просякнутий солодкістю, так усе це, вияв Я, цілком просякнуте мною.

[2.7] З незнання себе постає ілюзія. У самопізнанні ілюзія згасає. Мотузка --- не змія, хоч і може нею здатися.

[2.8] Світло --- моя справжня природа, ні менше, ні більше за це. Всесвіт проявляється на моєму видноколі.

[2.9] Ілюзія незнання постає в мені так, як змія постає в мотузці, як міраж води на сонячному обрії і як срібло в перламутрі.
